\section{Conclusion}
\label{sec:conclusion}

We conducted a controlled empirical study of capability-level diversity in
multi-agent LLM systems using a custom 30-question adversarial reasoning benchmark
across five cognitive trap categories.
We evaluated five conditions spanning a diversity continuum---single agents,
homogeneous three-agent ensembles, heterogeneous two-small-plus-one-large ensembles,
and a structured debate protocol---using real API calls to Anthropic's
\haiku and \sonnet models.

Our main finding is that within-family capability diversity produces measurably
lower error correlation ($r{=}0.464$ vs.\ $r{=}0.695$ for within-capability pairs),
confirming the theoretical diversity hypothesis at the correlation level.
However, this advantage does not translate to statistically significant accuracy gains
on our benchmark, where both models individually achieve 93.3\% accuracy.
The debate protocol degrades performance due to sycophantic revision,
and the dominant shared failure mode---a safety-aligned medical refusal---is
orthogonal to capability-level diversity and cannot be overcome by any within-family ensemble.

Taken together, our results support \citeauthor{rosales2025diverse}'s~\citeyearpar{rosales2025diverse}
argument that complementary failure modes, not mere diversity, are what matter for ensemble benefit.
Within-family models share not just capabilities but safety alignment properties,
creating systematic shared biases that transcend capability-level diversity.

\para{Key takeaway.}
Capability-level diversity within a single LLM family reduces error correlation
but does not reliably improve accuracy because shared alignment properties create
correlated failures that cannot be overcome by mixing model sizes.

\para{Future work.}
Three directions follow naturally from our findings.
First, a \emph{cross-family study} using OpenAI, Anthropic, and open-source models
(GPT-4o, Claude Sonnet, LLaMA-3) would test whether cross-family diversity
yields complementary failure modes even on safety-sensitive questions.
Second, a \emph{larger and novel adversarial benchmark} (200+ questions,
programmatically generated, not present in training data) would provide
the statistical power needed to detect the diversity effects that our
$n{=}30$ near-saturated benchmark cannot resolve.
Third, \emph{cascade architectures} that route questions by difficulty
rather than aggregating all agents on all questions offer a promising
alternative to ensemble voting, trading diversity benefits for cost efficiency
while avoiding the sycophancy pitfalls of debate.
